\scanpage{90}
Ce groupe est engendré par ${\unsure{G|\mathfrak S|\mathcal S}}_3$ et $6$ éléments $\rho_\sigma$ ($\sigma\in G_3$) \deleted{$\rho_\sigma=\sigma$} avec relations
\[
\begin{aligned}
\left\{\begin{aligned}
\rho_\sigma&=\sigma\rho_1\sigma^{-1}\quad\inserted{(\sigma\in G_3)}
  \quad[\text{donc }\rho_{\tau\sigma}=\tau\rho_\sigma\tau^{-1}\ \inserted{(\sigma,\tau\in G_3)}]\\
\rho_1^2&=1\qquad [\Rightarrow \rho_\sigma^2=1]\\
\rho_1\rho_\pi&=\rho_\pi\rho_1\qquad(\pi=\pi_{01}\ \text{transposition de }0,1)
\end{aligned}\right.\\
[\Rightarrow(\rho_\sigma\rho_{\sigma\pi}=\rho_{\sigma\pi}\rho_\sigma\qquad]
\end{aligned}
\]
Le groupe \deleted{\unsure{qui sera}} (des ``opérations'' élémentaires discrètes'') opère p.ex. sur l'ens des $72+24=96$ syst. de Schwartz \inserted{$\mathrm{Sch}(H)$} d'un groupe tétraédral $H(\simeq A_4)$, donc sur l'ens $\widetilde{\mathrm{Sch}}(H)$ formé de $8$ éléments $[(3,3,2)^+\ (3,3,2)^-$ et les transformés par $G_3^+$, $(3,3,3)^+$ et $(3,3,3)^-]$, donc sur l'ens à $4$ élts, des classes sous \deleted{Comme on a explicité \unsure{déj.}} $\widetilde H=\mathrm{Aut}(H)$ $(\simeq G_{\unsure{4|3}})$, formé de $(3,3,2)$ $(3,\deleted{\unsure{3}}2,3)$ $(2,3,3)$ et de $(3,3,3)$. Comme on a explicité l'action de $G_3$ sur ces ens, il \unsure{restera|resterait} à expliciter l'action de $\unsure{E|\mathcal E}_0$, ou ce qui revient dès lors au $=$, de $\rho_1$ i.e. $\rho_{-1,1.}$ L'exercice est immédiat en fonction de ce qui a été vu. \deleted{\unsure{pour}} Il serait intéressant d'expliciter les gpes quotients finis de $G_3\cdot E_0$ qui opèrent \uline{fidmt} sur les \unsure{$3/$ens|$3$ ens.} précédents \ldots
% archive sheet number: 89
